\section{统计损伤证书与有限样本应用}
\label{app:statistical}
本节给出主定理记录关系的一项应用。统计实验比较\cite{blackwell1953}和分布距离不等式\cite{sason2016}
用于描述记录的信息与期望约束，分布约束优化\cite{bertsimas2005}提供相应求界背景；
随机单比特编码\cite{benbasat2021,suresh2017}则限制每次输出的位数。
以下构造采用正文已定义的 $P=P_k^+$、$Y=\ind\{\gamma_{k+1}=0\}$ 以及常数 $M,q,\Omega_k$。
先给出固定核的统计损伤证书，再说明其独立样本保证和数值代价。
\subsection{固定记录核与统计损伤证书}
式\eqref{eq:witness}中的加权静默事件反映条件估计偏差。
将其与协方差迹项合并，可构造取值在 $[0,1]$ 的实值得分，
使得分均值之差直接控制真实均方误差增量。
令 $p_{\max}=\max_{P\in\mathcal P_k^+}\operatorname{tr}P$，定义
\begin{equation}
 \begin{aligned}
 L&=M+q\Omega_k,\\
 T(P,Y)&=\frac{w(P)Y+q(p_{\max}-\operatorname{tr}P)}{L},\\
 B\mid(P,Y)&\sim\operatorname{Bernoulli}(T(P,Y)).
 \end{aligned}
 \label{eq:score}
\end{equation}
二值化使用与系统及攻击相独立的额外随机数。
无攻击与受攻运行采用同一固定记录核；由于 $0\le T\le1$，
实值得分与二值输出均可通过样本均值估计其期望。
\begin{proposition}[向量系统的统计损伤证书]
\label{prop:certificate}
对主定理的全部允许策略及 $b\le b_{\rm ref}$，有
\begin{equation}
 \begin{aligned}
 J_k&\le C_B(\E_0T-\E_\pi T)
       =C_B(\E_0B-\E_\pi B),\\
 C_B&=L/q=C_{\rm record}.
 \end{aligned}
 \label{eq:signed-certificate}
\end{equation}
\end{proposition}
\begin{proof}
由式\eqref{eq:score}，
\begin{equation}
 \begin{aligned}
 L(\E_0T-\E_\pi T)
 &=1-\E_\pi e^{-D^\top H(P)D}\\
 &\quad+q(\E_\pi\operatorname{tr}P-\E_0\operatorname{tr}P)\\
 &\ge qJ_k.
 \end{aligned}
 \label{eq:score-identity}
\end{equation}
条件期望又给出 $\E B=\E T$，故结论成立。
\end{proof}
该构造给出适用于 $n$ 维系统的可行证书。
常数沿用主定理中的统一半径界与最小特征值，保守性也由此传递到证书中；
其在全部记录核中的最优性尚未建立。

\subsection{记录成本与二值化方差}
矩阵标签 $P_k^+$ 有 $n(n+1)/2$ 个对称分量；作为有限事件历史的函数，
它也可以用模型依赖的标签索引表示。
当历史到精确标签的映射为单射时，标签与历史具有相同的统计区分信息。
此时，精确矩阵标签提供的是事件历史的等价表示。
标签量化会合并不同历史，对量化记录的损伤保证需要重新推导，
以计入精确 $w(P)$ 被替换后产生的误差。

二值记录 $B$ 每次输出一位，其生成需要 $P$、下一时刻事件及预计算的记录核。
因此，输出位数、在线协方差存储与更新、离线历史枚举及核计算是不同的成本项。
实值得分 $T$ 保留相同均值，而全方差公式给出
\begin{equation}
 \operatorname{Var}(B)=\operatorname{Var}(T)+\E[T(1-T)].
 \label{eq:variance}
\end{equation}
在独立等样本数条件下，随机二值化保持均值，并增加非负的抽样方差项。
二值输出的传输或存储收益应与这一统计代价共同评估。

\subsection{独立实验下的有限样本置信上界}
\label{sec:finite-sample}

命题\ref{prop:certificate}用总体期望差控制均方误差增量。
从有限记录计算损伤上界时，还需计入均值估计的抽样误差。
固定记录核、攻击策略分布及正整数样本量，取 $0<\alpha<1$，记
$p_0=\E_0B$、$p_a=\E_\pi B$。式\eqref{eq:signed-certificate}给出
有符号关系 $J_k\le C_B(p_0-p_a)$，因此只需构造概率差的单侧上界。
以下实验均从相同初始分布独立重置，攻击私有随机数
与记录随机化也在每次实验中独立生成。

若无攻击时的概率 $p_0$ 已知，以 $n_a$ 次受攻实验得到均值 $\hat p_a$，
Hoeffding不等式\cite{hoeffding1963}给出：以至少 $1-\alpha$ 的概率，
\begin{equation}
 J_k\le C_B\left[p_0-\hat p_a+
 \sqrt{\frac{\log(1/\alpha)}{2n_a}}\right].
 \label{eq:fs-known}
\end{equation}
若另以独立的 $n_0$ 次无攻击实验估计 $p_0$，则
\begin{equation}
 J_k\le C_B\left[\hat p_0-\hat p_a+
 \sqrt{\frac{\log(1/\alpha)}{2}
 \left(\frac1{n_0}+\frac1{n_a}\right)}\right]
 \label{eq:fs-two}
\end{equation}
同样具有至少 $1-\alpha$ 的置信水平。
为证明式\eqref{eq:fs-two}，令 $X=\hat p_0-\hat p_a$。其中无攻击样本的权重为
$1/n_0$，受攻样本的权重为 $-1/n_a$，各独立加权项的取值区间长度
平方之和为 $1/n_0+1/n_a$。直接应用Hoeffding下尾界即得
$\Pr\{\E X-X>t\}\le\exp[-2t^2/(1/n_0+1/n_a)]$。
该推导直接控制两组样本均值之差，避免了分别构造置信区间再相加所带来的额外余量。

上述公式也适用于 $[0,1]$ 内的实值得分 $T$，将概率估计替换为得分均值即可。
Hoeffding界仅使用取值区间，因此两种输出具有相同形式的置信余量，
式\eqref{eq:variance}中的方差差异不影响该余量。
实际计算以观测均值代入，并将所得上界与预定允许损伤比较；
若上界仍高于该阈值，则本次数据不足以完成认证。
这里的认证对象是给定策略分布下的期望损伤。

\textbf{给定损伤阈值的样本量规划。}
令允许损伤为 $\tau$，总体概率差为 $\delta=p_0-p_a$，并记
$m=\tau/C_B-\delta$。在参考概率及记录核均精确的条件下，若 $m>0$，令经验概率差
等于总体差并反解式\eqref{eq:fs-known}、\eqref{eq:fs-two}，得到
\begin{equation}
 \begin{aligned}
 n_a&\ge\left\lceil\frac{\log(1/\alpha)}{2m^2}\right\rceil
 &&(p_0\text{已知}),\\
 n_0=n_a&\ge\left\lceil\frac{\log(1/\alpha)}{m^2}\right\rceil
 &&(\text{两组等量采样}).
 \end{aligned}
 \label{eq:fs-planning}
\end{equation}
式\eqref{eq:fs-planning}给出经验概率差恰等于总体差时的样本量规划值，
用于衡量置信余量的规模。它既非必要样本量，也不保证以 $1-\alpha$ 概率成功认证，
因为实际经验均值仍存在随机偏差。
后文以二维工况给出相应数值；认证时仍须代入实际样本均值。
若 $m\le0$，则即使经验概率差等于总体差，有限样本上界也无法达到目标。

\textbf{实现误差与适用条件。}
上述结论以模型证书成立和记录核准确实现为前提。若使用的无攻击参考值
$\tilde p_0$ 满足 $|\tilde p_0-p_0|\le\delta_0$，则在式
\eqref{eq:fs-known}中以 $\tilde p_0$ 替换 $p_0$，并在括号内加上 $\delta_0$。
若实际二值输出为 $\widetilde B$，且已验证
$|\E_i\widetilde B-\E_iB|\le\eta_i$（$i=0,a$，其中 $\E_a=\E_\pi$），
则以实际输出的样本均值代入式\eqref{eq:fs-two}，并在括号内加上 $\eta_0+\eta_a$。
无攻击参考采用规定核的理论概率时，式\eqref{eq:fs-known}的括号内只需加 $\eta_a$。
这些概率误差均需乘以外部系数 $C_B$ 才转换为损伤余量；
模型失配及证书常数的有效性仍需单独核验。
记录核、样本数和置信水平应事先确定；若由数据选择记录核，
则应另用独立样本验证。以上置信保证适用于独立重置实验，
对于单条连续轨迹中的相关记录，需要另行建立相应的统计界。

\subsection{二维算例中的统计代价}
采用 $b_{\rm ref}=0.2$ 预先固定记录核，并在 $b=0.1$ 的两拍符号攻击下评价该核。
首先比较同一输入记录产生的实值得分 $T$ 与二值输出 $B$，结果见表\ref{tab:vector-variance}。
两者的总体均值相同，二值输出的方差则分别增加0.014176和0.043915（各列独立舍入），
与 $\E[T(1-T)]$ 一致。因此，将输出限制为单比特会在相同实验次数下引入额外的抽样方差。

\begin{table}[H]
\centering\small
\tabcaption{同攻击下实值得分与二值输出的方差}{Variances of the real-valued score and binary output under the same attack}
\label{tab:vector-variance}
\begin{tabular}{lrrr}
\toprule
工况 & $\operatorname{Var}(T)$ & $\operatorname{Var}(B)$ & 新增方差\\
\midrule
S & 0.222770 & 0.236945 & 0.014176 \\
U & 0.203677 & 0.247592 & 0.043915 \\
\bottomrule
\end{tabular}
\end{table}

\begin{table}[H]
\centering\small
\tabcaption{总体中心处的95\%损伤上界示例}{Illustrative 95\% damage bounds evaluated at the population means}
\label{tab:vector-confidence}
\begin{tabular}{lrrrr}
\toprule
工况 & $C_B$ & 统计损伤证书 & 参考已知 & 两组估计\\
\midrule
S & 6.884 & 0.000458 & 0.085 & 0.120 \\
U & 5.869 & 0.000902 & 0.073 & 0.102 \\
\bottomrule
\end{tabular}
\end{table}
表\ref{tab:vector-confidence}使用每组 $10^4$ 次独立实验对应的单侧Hoeffding余量，
并令经验均值等于总体均值；这是置信界公式的示意计算，并非新增采样结果。
S、U的统计损伤证书分别约为 $0.000458$ 和 $0.000902$，均高于相应真实损伤；
加入有限样本余量后，上界明显增大。
若目标阈值设为各自统计损伤证书加 $0.01$，
式\eqref{eq:fs-planning}的已知参考中心规划值分别为709866和515964次；
两组等量估计时分别为每组1419732和1031927次。
这些总体中心处的规划值由本文的证书常数和Hoeffding界确定，
并非统计认证所必需的最小实验次数。输出位数与实验次数应分别计入成本。

最后，以仅利用峰值约束的上界为基准，评价记录提供的额外信息。
当 $b=0.1$ 时，式\eqref{eq:peak-bound}给出的S、U损伤上界分别为0.033677和0.132695。
在表\ref{tab:vector-confidence}的总体中心条件下，每组 $10^4$ 个样本所得的两种统计界
均未收紧S的幅值界；对于U，无攻击参考已知和两组估计时则分别收紧约45.2\%和22.8\%。
由未舍入参数反解单侧余量可得：若经验均值等于总体均值，
使该Hoeffding统计界严格小于本文幅值界的最小整数样本数，
在S中为64326次受攻实验，或两组各128652次；在U中相应为2971次，或两组各5942次。
这些阈值对应总体中心处的界限交点，实际认证结果仍随样本均值而变化。
记录能否收紧本文的幅值界取决于系统参数和样本量。
实际使用时，可将有限样本置信上界与 $C_{\rm peak}b^2$ 取较小者；
表\ref{tab:vector-confidence}单列统计界，以区分统计损伤证书与采样余量。

